\worldchapter{\trans{weapon_modifications_title}}{appendix-weaponmodifications}
\label{ch:appendix-weaponmodifications}

\begin{multicols}{2}

    \section{Visiere}



    \subsection*{Armbrust-Eisen-Visier}

    \begin{pexample}
    Das Anbringen eines eisernen Visiers an der Armbrust erhöht die Treffsicherheit und den möglichen Schaden bei einem Treffer.
    \end{pexample}


    \begin{tabular}{@{}ll@{}}
        \textbf{\trans{price_label}:} & 600 \\
        \textbf{Genauigkeit} & +1\\
        \textbf{Schadenspotential} & +1\\
    \end{tabular}

    \wbif{phasesix}{
        \vspace{3mm}
\faIcon{cube}\hspace{0.5em}    }{}


    \section{Sondervorrichtung}



    \subsection*{Schnellziehköcher}

    \begin{pexample}
    Dieser Köcher ist so konstruiert, dass ein Pfeil viel schneller auf die Sehne eines Bogens gesetzt werden kann.
    \end{pexample}


    \begin{tabular}{@{}ll@{}}
        \textbf{\trans{price_label}:} & 200 \\
        \textbf{Nachladeaktionen} & -1\\
    \end{tabular}

    \wbif{phasesix}{
        \vspace{3mm}
\faIcon{cube}\hspace{0.5em}    }{}



    \subsection*{Schnellziehschlinge}

    \begin{pexample}
    Diese Vorrichtung an der Waffe ermöglicht ein schnelles Ziehen und Schießen.
    \end{pexample}


    \begin{tabular}{@{}ll@{}}
        \textbf{\trans{price_label}:} & 200 \\
        \textbf{Vorbereitung} & -1\\
    \end{tabular}

    \wbif{phasesix}{
        \vspace{3mm}
\faIcon{cube}\hspace{0.5em}    }{}



    \subsection*{Gesegnet}

    \begin{pexample}
    Die Waffe wurde von einem Priester geweiht. Auf ihr liegt der Segen einer höheren Wesenheit, sie hat besondere Fähigkeiten und wirkt stärker gegen die Mächte des Bösen.
    \end{pexample}

    Ergebnisse des Trefferwurfes, die eine 1 ergeben, dürfen einmal wiederholt werden. Die Anzahl der Treffer gegen Dämonen und Geister wird verdoppelt.

    \begin{tabular}{@{}ll@{}}
        \textbf{\trans{price_label}:} & 500 \\
        \textbf{Schadenspotential} & +1\\
    \end{tabular}

    \wbif{phasesix}{
        \vspace{3mm}
\faIcon{ankh}\hspace{0.5em}\faIcon{spider}\hspace{0.5em}    }{}


    \section{Griffe}



    \subsection*{Lederumwickelter Griff}

    \begin{pexample}
    Ein mit Leder überzogener Griff verbessert die Handhabung der Waffe und erhöht das Schadenspotential.
    \end{pexample}


    \begin{tabular}{@{}ll@{}}
        \textbf{\trans{price_label}:} & 80 \\
        \textbf{Schadenspotential} & +1\\
    \end{tabular}

    \wbif{phasesix}{
        \vspace{3mm}
\faIcon{eye}\hspace{0.5em}\faIcon{chess-rook}\hspace{0.5em}\faIcon{umbrella}\hspace{0.5em}\faIcon{flag}\hspace{0.5em}\faIcon{plane}\hspace{0.5em}\faIcon{microchip}\hspace{0.5em}    }{}



    \subsection*{Hartholzgriff}

    \begin{pexample}
    Ein Hartholzgriff verbessert die Handhabung der Waffe und erhöht das Schadenspotential und die Präzision.
    \end{pexample}


    \begin{tabular}{@{}ll@{}}
        \textbf{\trans{price_label}:} & 200 \\
        \textbf{Genauigkeit} & +1\\
        \textbf{Schadenspotential} & +1\\
    \end{tabular}

    \wbif{phasesix}{
        \vspace{3mm}
\faIcon{cube}\hspace{0.5em}    }{}


    \section{Munition}



    \subsection*{Flintenlaufgeschosse}

    \begin{pexample}
    Diese Munition ermöglicht es, aus einem Gewehr ein einzelnes Geschoss abzufeuern, das mehr Schaden anrichtet und die Reichweite des Gewehrs erhöht.
    \end{pexample}


    \begin{tabular}{@{}ll@{}}
        \textbf{\trans{price_label}:} & 100 \\
        \textbf{Schadenspotential} & +1\\
        \textbf{Reichweite} & +5\\
    \end{tabular}

    \wbif{phasesix}{
        \vspace{3mm}
\faIcon{cube}\hspace{0.5em}    }{}



    \subsection*{Erweitertes Magazin (Pistolen)}

    \begin{pexample}
    Das erweiterte Magazin fasst 7 Schuss mehr und kann für Pistolen verwendet werden.
    \end{pexample}


    \begin{tabular}{@{}ll@{}}
        \textbf{\trans{price_label}:} & 80 \\
        \textbf{Kapaziät} & +7\\
    \end{tabular}

    \wbif{phasesix}{
        \vspace{3mm}
\faIcon{cube}\hspace{0.5em}    }{}



    \subsection*{Erweitrertes Magazin (Maschinengewehre)}

    \begin{pexample}
    Das erweiterte Magazin fasst 20 Schuss mehr und kann für Maschinengewehre verwendet werden.
    \end{pexample}


    \begin{tabular}{@{}ll@{}}
        \textbf{\trans{price_label}:} & 150 \\
        \textbf{Kapaziät} & +20\\
    \end{tabular}

    \wbif{phasesix}{
        \vspace{3mm}
\faIcon{cube}\hspace{0.5em}    }{}



    \subsection*{Kieselsteine}

    \begin{pexample}
    Einfache Kieselsteine zur Benutzung in einer Schleuder oder Zwille.
    \end{pexample}


    \begin{tabular}{@{}ll@{}}
        \textbf{\trans{price_label}:} & 2 \\
    \end{tabular}

    \wbif{phasesix}{
        \vspace{3mm}
\faIcon{cube}\hspace{0.5em}    }{}



    \subsection*{Eisenkugeln}

    \begin{pexample}
    Eisenkugeln fügen dem Gegner mehr Schaden zu, wenn sie anstelle von Steinen in einer Schleuder verwendet werden.
    \end{pexample}


    \begin{tabular}{@{}ll@{}}
        \textbf{\trans{price_label}:} & 10 \\
        \textbf{Schadenspotential} & +2\\
    \end{tabular}

    \wbif{phasesix}{
        \vspace{3mm}
\faIcon{cube}\hspace{0.5em}    }{}



    \subsection*{Giftpfeile}

    \begin{pexample}
    Giftpfeile haben eine spezielle Spitze, an der das Gift besonders gut haftet. Diese Pfeile verursachen Vergiftungen, die der Stärke des verwendeten Giftes entsprechen.
    \end{pexample}


    \begin{tabular}{@{}ll@{}}
        \textbf{\trans{price_label}:} & 20 \\
        \textbf{Giftkanal} & +1\\
    \end{tabular}

    \wbif{phasesix}{
        \vspace{3mm}
\faIcon{cube}\hspace{0.5em}    }{}



    \subsection*{Explosive Pfeile}

    \begin{pexample}
    Durch eine spezielle Vorrichtung an der Pfeilspitze explodieren die Pfeile beim Aufprall.
    \end{pexample}


    \begin{tabular}{@{}ll@{}}
        \textbf{\trans{price_label}:} & 700 \\
        \textbf{Flächenschaden} & +2\\
    \end{tabular}

    \wbif{phasesix}{
        \vspace{3mm}
\faIcon{cube}\hspace{0.5em}    }{}


    \section{Klinge}



    \subsection*{Aufgerauhte Klinge}

    \begin{pexample}
    Wird die Klinge einer Waffe aufgerauht, so verringert sich zwar die Durchschlagskraft der Waffe, ein Schlag verursacht jedoch stark blutende Wunden.
    \end{pexample}


    \begin{tabular}{@{}ll@{}}
        \textbf{\trans{price_label}:} & 100 \\
        \textbf{Durchschlag} & -1\\
        \textbf{Blutend} & +2\\
    \end{tabular}

    \wbif{phasesix}{
        \vspace{3mm}
\faIcon{cube}\hspace{0.5em}    }{}



    \subsection*{Gehärtete Klinge}

    \begin{pexample}
    Die gehärtete Klinge erhöht die Durchschlagskraft und das Schadenspotential der Waffe.
    \end{pexample}


    \begin{tabular}{@{}ll@{}}
        \textbf{\trans{price_label}:} & 200 \\
        \textbf{Schadenspotential} & +1\\
        \textbf{Durchschlag} & +1\\
    \end{tabular}

    \wbif{phasesix}{
        \vspace{3mm}
\faIcon{cube}\hspace{0.5em}    }{}



    \subsection*{Gravierte Klinge}

    \begin{pexample}
    Die Klinge der Waffe enthält eine besondere Gravur.
    \end{pexample}


    \begin{tabular}{@{}ll@{}}
        \textbf{\trans{price_label}:} & 100 \\
        \textbf{Schadenspotential} & +1\\
    \end{tabular}

    \wbif{phasesix}{
        \vspace{3mm}
\faIcon{cube}\hspace{0.5em}    }{}



    \subsection*{Giftkanal}

    \begin{pexample}
    Eine Kerbe zum Auftragen von Gift. Klingenwaffen können damit modifiziert werden. Vergiftet in der Stärke des verwendeten Giftes.
    \end{pexample}


    \begin{tabular}{@{}ll@{}}
        \textbf{\trans{price_label}:} & 250 \\
        \textbf{Giftkanal} & +1\\
    \end{tabular}

    \wbif{phasesix}{
        \vspace{3mm}
\faIcon{cube}\hspace{0.5em}    }{}



    \subsection*{Gekrümmte Klinge}

    \begin{pexample}
    Wenn die Waffe eine gekrümmte Klinge hat, wird die Reichweite erhöht und die Wunde blutet, da die Waffe häufiger ungeschützte Körperteile erreicht.
Eine bestehende Waffe kann nicht durch einen Schmied auf eine gekrümmte Klinge umgerüstet werden, dies muss bei neuen Waffen direkt beauftragt werden.
    \end{pexample}


    \begin{tabular}{@{}ll@{}}
        \textbf{\trans{price_label}:} & 300 \\
        \textbf{Reichweite} & +1\\
        \textbf{Blutend} & +1\\
    \end{tabular}

    \wbif{phasesix}{
        \vspace{3mm}
\faIcon{cube}\hspace{0.5em}    }{}



    \subsection*{Gezackte Schneide}

    \begin{pexample}
    Eine gezackte Schneide verursacht stark blutende Wunden.
    \end{pexample}


    \begin{tabular}{@{}ll@{}}
        \textbf{\trans{price_label}:} & 400 \\
        \textbf{Blutend} & +1\\
    \end{tabular}

    \wbif{phasesix}{
        \vspace{3mm}
\faIcon{cube}\hspace{0.5em}    }{}



    \subsection*{Verzauberung}

    \begin{pexample}
    Eine Verzauberung durch einen gesungenen Spruch
    \end{pexample}

    Die Waffe erhält +2 Schadenspotential

    \begin{tabular}{@{}ll@{}}
        \textbf{\trans{price_label}:} & 400 \\
        \textbf{Schadenspotential} & +2\\
    \end{tabular}

    \wbif{phasesix}{
        \vspace{3mm}
\faIcon{magic}\hspace{0.5em}    }{}


\end{multicols}